\chapter{外部数据及其约束范围}
\label{ch:data}

公开数据只约束那些它真正测到的量。一份等温线不能验证冷点估计，一份热导率表不能验证探针的信息增益。本章记录已经取得并读过的数据，以及尚缺、因而不能写入主结果的数据。

\section{已读入的等温线}

Zenodo 记录 \texttt{10.5281/zenodo.21671925}，许可 CC BY 4.0。内容是从已发表文献整理的食品水分吸附/解吸曲线，以及经典 GAB 与重参数 GAB 的拟合。读入的是拟合表，共 200 行曲线记录；数据说明中的汇总数是 192 条保留曲线。清洁拟合（诊断码 0）140 条。

在清洁拟合上，$X_m$ 的十分位、中位、九十分位约为 $0.037$、$0.074$、$0.144$；$C$ 约为 $4$、$11$、$45$；$K$ 约为 $0.61$、$0.77$、$0.96$。这些分位只描述该表的清洁子集，不描述全部食品。

该表的食品类别以谷物、果蔬和粉体为主。按名称检索，肉类曲线条数为 0。土豆有 6 条清洁拟合，温度只到 $60^\circ\mathrm{C}$：一组的 $X_m$ 从 $30^\circ\mathrm{C}$ 的 $0.057$ 降到 $60^\circ\mathrm{C}$ 的 $0.036$，$K$ 从 $0.91$ 降到 $0.76$；另一组 $X_m$ 从 $0.076$ 降到 $0.052$。方向与吸附热随温度变化一致，但 $60^\circ\mathrm{C}$ 仍低于烤制表面温度。

\begin{table}[H]
\centering
\caption{这份等温线能够约束和不能约束的对象}
\begin{tabularx}{\textwidth}{lX}
\toprule
能 & 不能 \\
\midrule
淀粉类 $a_w(M,T)$ 在测定温度内的先验区间 & 肉类、脂肪组织的等温线 \\
GAB 三参数的数量级与相关结构 & 蒸烤过程中的动态非平衡含水 \\
温度升高时 $X_m$ 下降这一方向 & $100^\circ\mathrm{C}$ 以上的外推数值 \\
灵敏度扫描的参数上下界 & 冷点、探针信息增益、控制效果 \\
\bottomrule
\end{tabularx}
\end{table}

使用时保留记录号、类别、温度、吸附或解吸、诊断码和参数标准误。诊断码不为 0 的行不进入先验。成稿引用该数据集的 DOI 与许可，不把拟合值写成新的实验测定。

\section{尚缺的测量与仿真如何补位}

下列对象在 Zenodo 与 Dryad 的数据集检索中没有得到可引用的测量文件。检索词覆盖肉类导热、烤制中心温度、面包烘烤和微生物热失活。返回结果是晶体导热、酵母烘烤结构、色谱或无关材料，没有进入参数表。

\begin{table}[H]
\centering
\caption{缺测对象的临时约束}
\begin{tabularx}{\textwidth}{lXl}
\toprule
对象 & 在测量到达前 & 到达后替换成 \\
\midrule
肌肉与脂肪导热 & 文献区间加 S1 的灵敏度 & 该文件的温度—含水表 \\
内部温度与质量历程 & 只由 T1/T2 互相对照 & 外部残差核对 \\
腔体换热与混合 & T3 边界库，标记为计算边界 & 实测 $h(v)$ \\
反应进度 & 动力学文献的数量级，S2 扫描 $E_a$ & 同步测量的进度 \\
微生物绝对初始量 & 不进入状态 & 仍不进入状态 \\
\bottomrule
\end{tabularx}
\end{table}

仿真补的是覆盖范围，不是测量。S1 的 1\,620 条轨迹不能写成“相当于 1\,620 次实验”。参数文件里来源字段只允许三个值：\texttt{doi:10.5281/zenodo.21671925}、\texttt{literature-interval}、\texttt{simulation}。第三个值不能出现在摘要的验证句里。

后续测量文件一旦取得，先按第 \ref{ch:data} 章的入口做核对，再决定是否缩小 S1 里该参数的抽样区间。不重跑全部格子，除非核对残差改变了策略排序。

\section{数据进入模型的唯一入口}

外部测量不直接覆盖在线后验。入口只有两处。

\begin{enumerate}
\item \textbf{先验区间。} 清洁 GAB 参数写成类别分布的上下界。在线更新不得跑出该界，除非本次质量与湿度残差证明界本身与样品矛盾，并在日志中标记。
\item \textbf{外部核对。} 若未来的公开温度历程与仿真边界足够接近，比较的是 $T(t)$ 与质量损失的残差。残差过大时收缩几何或边界的适用范围，不改命题 \ref{prop:steady}--\ref{prop:shield} 的陈述。
\end{enumerate}

演示下载的参数文件要带数据来源字段：哪些常数来自该 DOI，哪些仍是无量纲例证。无来源的肉类等温线不得标成来自这份数据。
